\documentclass{article}
\usepackage{amsmath,amssymb}
\usepackage{xurl}
\usepackage[colorlinks=true,linkcolor=blue,citecolor=blue,urlcolor=blue]{hyperref}
% Shared commands for notes in docs/paper-gaps/.

\DeclareUnicodeCharacter{2080}{\ifmmode{}_{0}\else\textsubscript{0}\fi}
\DeclareUnicodeCharacter{2081}{\ifmmode{}_{1}\else\textsubscript{1}\fi}
\DeclareUnicodeCharacter{2082}{\ifmmode{}_{2}\else\textsubscript{2}\fi}
\DeclareUnicodeCharacter{2099}{\ifmmode{}_{n}\else\textsubscript{n}\fi}

\newcommand{\Lean}{\textsc{Lean}}
\ExplSyntaxOn
\NewDocumentCommand{\leanid}{m}{
  \ifmmode
    \textsf{\detokenize{#1}}
  \else
    \group_begin:
    \urlstyle{sf}
    \tl_set:Nn \l_tmpa_tl {#1}
    \tl_replace_all:Nnn \l_tmpa_tl {\_} {_}
    \exp_args:NV \path \l_tmpa_tl
    \group_end:
  \fi
}
\ExplSyntaxOff
\providecommand{\tr}{\operatorname{tr}}
\newcommand{\tnleanIssueBase}{https://github.com/LionSR/TNLean/issues/}
\newcommand{\tnleanPrBase}{https://github.com/LionSR/TNLean/pull/}
\newcommand{\tnleanDocBase}{https://sirui-lu.com/TNLean/docs/}
\newcommand{\ghissue}[1]{\href{\tnleanIssueBase#1}{GitHub issue \##1}}
\newcommand{\ghpr}[1]{\href{\tnleanPrBase#1}{PR \##1}}
\newcommand{\doclink}[1]{\href{\tnleanDocBase#1.html}{\path{#1}}}

% Dirac notation and the identity operator, as in blueprint/src/macros/common.tex.
% \providecommand keeps note-local definitions authoritative.
\usepackage{dsfont}
\providecommand{\ket}[1]{|#1\rangle}
\providecommand{\bra}[1]{\langle #1|}
\providecommand{\braket}[2]{\langle #1 | #2 \rangle}
\providecommand{\ketbra}[2]{|#1\rangle\!\langle #2|}
\providecommand{\Id}{\mathds{1}}
\providecommand{\one}{\mathds{1}}
\providecommand{\C}{\mathbb{C}}
\providecommand{\MN}[1]{M_{#1}(\C)}
\providecommand{\MD}{M_D(\C)}

% Verdict marker: \gapnote{<kind>}{<status>}. Kind is the policy
% classification (clarification, local-correction, scope-restriction,
% unfaithful, false-source, open-gap); status is open, wip (elimination
% underway), resolved, or historical. Machine-read by the site index;
% prints a verdict line.
\newcommand{\gapnote}[2]{\par\noindent\textbf{Verdict:} #1 (#2).\par}

\title{The Constants of the Buffered Discount in the One-Dimensional Area Law}
\author{}
\date{October 16, 2026}
\begin{document}
\maketitle
\gapnote{local-correction}{open}

\section{The source statement}

Lemma~4.2 of the one-dimensional note~\cite{TNLean2026OneDimensional}
(label \texttt{lem:buffer}, file \texttt{part1\_arealaw.tex}, lines~360--364)
is the buffered discount. For an interval $X$ of the chain with the shells
$X_j=X^{+jR}$, $1\le j\le m$, the flanks $T_j=X_j\setminus X$ and the uniform
weights $a_j=2/m$, the note fixes the constants (equation \texttt{eq:CE},
lines~321--328)
\begin{align}
    C_E
        &= 8\,N_R\,J\,q^{2(R+1)},
    \label{eq:tnlean_one_dimensional_discount_constants_source_energy_constant}\\
    m
        &= \Bigl\lceil\max\Bigl\{4,\ \frac{4C_E}{\Delta},\ 128\sqrt{(1+\vartheta)\mathcal B_R},\
        4096\,e\,(1+\vartheta)\,\mathcal B_R\Bigr\}\Bigr\rceil,
    \label{eq:tnlean_one_dimensional_discount_constants_source_shell_count}
\end{align}
with $\vartheta=J/\Delta$, notes $m\le C_1(q,R)(1+\vartheta)$
(equation \texttt{eq:mbound}, lines~331--333), and concludes
(equation \texttt{eq:discount})
\begin{align}
    S(X\,|\,T_m)
        &\le \tfrac12\,S(X)+\tfrac12+\frac{C_E}{\Delta}.
    \label{eq:tnlean_one_dimensional_discount_constants_source_discount}
\end{align}
In the proof, the energy of the filtered vector is bounded by
$\sum_j N_R\cdot 2a_j^2q^{2(R+1)}J=C_E/m$ (equation \texttt{eq:energy},
lines~458--463), so that by the second entry of $m$ it is at most $\Delta/4$
and the overlap defect $1-|\langle\Omega,\phi\rangle|^2\le C_E/(m\Delta)$ is
at most $1/4$ (equation \texttt{eq:overlap}, lines~465--467); the overlap
losses of the $m$ prefixes are then summed with $-\frac12\log(1-v)\le v$ to
$C_E/\Delta$, and the tail term
$512\,e\,\vartheta\,\mathcal B_R\sum_ja_j^2=2048\,e\,\vartheta\,\mathcal B_R/m$
is bounded by $\frac12$ through the fourth entry of $m$ (lines~505--516,
equation \texttt{eq:lowerN}). The closing step (equation \texttt{eq:final},
lines~606--608) then gives the interval bound
$C_{\mathrm{int}}=8mR\log q+1+2C_E/\Delta\le C_0(q,R)(1+J/\Delta)$ with
$C_0(q,R)=8C_1(q,R)R\log q+1+16N_Rq^{2(R+1)}$.

\section{What is proved and the deviations}

The formalization proves the discount on the chain for every interval
$X=[l,r)$ with the buffer $T=X^{+mR}\setminus X$, as an instance of the
finite-dimensional discount of the companion library,\footnote{Declarations:
\leanid{TNLean.MPS.AreaLaw.chainEnergyThreshold},
\leanid{TNLean.MPS.AreaLaw.chainTailThreshold},
\leanid{TNLean.MPS.AreaLaw.chainShellCount},
\leanid{TNLean.MPS.AreaLaw.chainDiscountConst},
\leanid{TNLean.MPS.AreaLaw.two_mul_condEntropy_interval_le}; the
finite-dimensional lemma is \leanid{Entropy.two_mul_condEntropy_le} of
QICLean.} in the form
\begin{align}
    2\bigl(S(X\cup T)-S(T)\bigr)
        &\le S(X)+\frac{4096\,e\,\vartheta\,\mathcal B_R}{m}+m\log2,
    \label{eq:tnlean_one_dimensional_discount_constants_formal_discount}\\
    m
        &= \Bigl\lceil\max\Bigl\{4,\ \frac{4C_E}{\Delta},\ 128\sqrt{(1+\vartheta)\mathcal B_R}\Bigr\}\Bigr\rceil .
    \label{eq:tnlean_one_dimensional_discount_constants_formal_shell_count}
\end{align}
The shells, the flanks, the weights, the crossing count $N_R$ and the cut
budget $\mathcal B_R$ are those of the note. The constants deviate from the
printed ones in three connected places, each at the level of an absolute
constant.

\begin{enumerate}
\item \emph{Energy threshold $\Delta/2$ in place of $\Delta/4$.} The
finite-dimensional lemma bounds the energy of the filtered vector by
$\sum_j4a_j^2\sum_{Y\text{ split by }X_j}(\dim\mathcal H_Y)^2J$, that is by
$16N_Rq^{2(R+1)}J/m=2C_E/m$, twice the printed estimate $C_E/m$. With the
same second entry $m\ge4C_E/\Delta$ of $m$ this is at most $\Delta/2$, so the
discount runs at the energy threshold $\Delta/2$, and the overlap defect of
every prefix is at most $1/2$ rather than $1/4$.
\item \emph{Overlap losses $m\log2$ in place of $2C_E/\Delta$.} The
finite-dimensional lemma keeps the overlap losses of the $m$ prefixes as
$-m\log(1-E/\Delta)$ with $E$ the energy threshold; at $E=\Delta/2$ this is
$m\log2$. The note instead uses $-\frac12\log(1-v)\le v$ at
$v=C_E/(m\Delta)\le\frac14$ and sums to $C_E/\Delta$ in $\log N$, that is
$2C_E/\Delta$ in the discount. Since $m\ge4C_E/\Delta$, the formal term
$m\log2$ dominates the printed $2C_E/\Delta$, and both are at most a
constant times $1+\vartheta$.
\item \emph{The tail term stays explicit and the fourth entry of $m$ is
omitted.} The error terms of the tail bound sum to
$\sum_j1024\,e\,\vartheta\,\mathcal B_ja_j^2\le4096\,e\,\vartheta\,\mathcal B_R/m$
in the discount, which is twice the printed $2048\,e\,\vartheta\,\mathcal B_R/m$
of $\log N$. The note bounds this by $1$ (that is, $\frac12$ in $\log N$)
through the fourth entry $m\ge4096\,e\,(1+\vartheta)\mathcal B_R$; the
formalization keeps the term $4096\,e\,\vartheta\,\mathcal B_R/m$ in the
discount constant and omits the fourth entry from $m$, so its $m$ is at most
the printed one.
\end{enumerate}

Consequently the formal interval bound of Theorem~A
is\footnote{Declarations:
\leanid{TNLean.MPS.AreaLaw.chainIntervalEntropyBound},
\leanid{TNLean.MPS.AreaLaw.chainShellConst},
\leanid{TNLean.MPS.AreaLaw.chainAreaLawConst},
\leanid{TNLean.MPS.AreaLaw.chainIntervalEntropyBound_le}.}
$C_{\mathrm{int}}=8mR\log q+4096\,e\,\vartheta\,\mathcal B_R/m+m\log2$ in
place of the printed $8mR\log q+1+2C_E/\Delta$, with the explicit
$C_1(q,R)=5+32N_Rq^{2(R+1)}+128\mathcal B_R$ in \texttt{eq:mbound} and
$C_0(q,R)=(8R\log q+\log2)\,C_1(q,R)+1024\,e\,\mathcal B_R$ in place of the
printed $C_0$. Every deviation is a change of absolute constants; the
structure of the argument, its hypotheses, and the linear dependence of the
area-law constant on $1+J/\Delta$ are those of the note.

\section{Elimination plan}

None of the deviations needs a formal change: the formal constants are valid
in place of the printed ones, and the printed ones are not sharper in any
respect that the note uses. A revision of the note may adopt the formal
constants (energy threshold $\Delta/2$, overlap losses $m\log2$, the tail
term kept explicit with the three-entry $m$), or the formalization may be
tightened to the printed ones by proving the finite-dimensional discount with
the sharper energy coefficient $2a_j^2$ and the inequality
$-\frac12\log(1-v)\le v$ for the overlap losses and by adding the fourth
entry to $m$. Until then the deviations are recorded on the module docstring
of \texttt{TNLean/MPS/AreaLaw/BufferDiscount.lean}, in the blueprint statements of
the buffered discount and of Theorem~A, which print the formal constants, and
in the status table of \path{docs/formalization/one-dimensional-chain.md}.

\bibliographystyle{alpha}
\bibliography{references}
\end{document}
